% Generated draft: requires mathematical and source review.
\documentclass[11pt]{article}
\usepackage[margin=25mm]{geometry}
\usepackage{fontspec}
\setmainfont{Latin Modern Roman}
\usepackage{amsmath,amssymb,mathtools,unicode-math}
\setmathfont{Latin Modern Math}
\usepackage{enumitem,needspace,xurl,hyperref}
\setlength{\parindent}{0pt}
\setlength{\parskip}{5pt}
\setlength{\emergencystretch}{3em}
\newcommand{\N}{\mathbb{N}}
\newcommand{\Z}{\mathbb{Z}}
\newcommand{\Q}{\mathbb{Q}}
\newcommand{\R}{\mathbb{R}}
\newcommand{\C}{\mathbb{C}}
\newcommand{\F}{\mathbb{F}}
\newcommand{\A}{\mathbb{A}}
\newcommand{\PP}{\mathbb{P}}
\newcommand{\E}{\mathbb{E}}
\newcommand{\eps}{\varepsilon}
\newcommand{\dd}{\,\mathrm{d}}
\begin{document}
\section*{UnsolvedMath ID: 800073}
\textbf{Source problem code:} catalogue entry
\subsection{Definitions and mathematical statement}
In the usual smooth interpretation, let \(M\) be a manifold of dimension \(4\), and let \(\mathcal F\) be a codimension-\(1\) foliation on \(M\). Its tangent bundle \(T\mathcal F\subset TM\) has rank \(3\). Let \(\xi\subset T\mathcal F\) be a rank-\(2\) subbundle that restricts to a contact structure on every leaf of \(\mathcal F\).

For smooth distributions \(E,F\subset TM\), write \([E,F]_x\) for the linear span of the values \([X,Y](x)\), where \(X\) and \(Y\) range over local smooth sections of \(E\) and \(F\), respectively. The leafwise contact condition is
\[
\xi_x+[\xi,\xi]_x=T_x\mathcal F
\qquad\text{for every }x\in M.
\]
Assume that there exist line subbundles \(L_1,L_2\subset\xi\) such that
\[
\xi=L_1\oplus L_2.
\]
A line subbundle here need not be oriented or trivial; the hypothesis does not itself supply global nonvanishing sections of the summands.

A smooth rank-\(2\) distribution \(D\subset TM\) is an Engel structure if its successive bracket distributions
\[
D^{(2)}_x=D_x+[D,D]_x,
\qquad
D^{(3)}_x=D^{(2)}_x+[D,D^{(2)}]_x
\]
satisfy
\[
\dim D^{(2)}_x=3,
\qquad
D^{(3)}_x=T_xM
\qquad\text{for every }x\in M.
\]

\textbf{Question (supplied Question 14).} Does every datum satisfying the preceding foliation, leafwise contact, and splitting hypotheses admit a perturbation to an Engel structure?

\textbf{Unresolved qualification.} The supplied formulation does not define the permitted perturbation. In particular, it does not specify a topology or degree of closeness, which components of the datum may change, whether the splitting is fixed or merely assumed to exist, or what relationship the resulting Engel structure must have to the original foliation. Consequently, the question above preserves the source wording but does not determine a unique quantified approximation statement. The source also does not explicitly specify regularity, compactness, boundary conditions, or orientability assumptions; no additional restrictions of those kinds are imposed here.
\subsection{Short English statement}
Given a \(4\)-manifold foliated by \(3\)-dimensional leaves, suppose a plane field is contact on each leaf and splits globally into two line subbundles. The question asks whether these data can be perturbed to an Engel distribution. The allowed perturbations and their required closeness are unspecified in the supplied record.
\subsection{Sources}
\begin{itemize}
\item Local UnsolvedMath registry: permanent ID \(800073\), exact source problem code \texttt{catalogue entry}, title \emph{Foliations and Diffeomorphism Groups --- Question 14}. The record attributes the problem to Nariman and Yazdi, with proposed year \(2025\), and labels it open. This status is supplied registry metadata, not an independently verified mathematical or literature assessment.
\item Nariman and Yazdi, \emph{Problem List on Foliations and Diffeomorphism Groups} (2025), Question 14. Supplied source URL: \url{https://arxiv.org/abs/2503.18274}. The supplied provenance records access on \texttt{2026-10-16} and extraction by \texttt{arxiv\_latex}; these are supplied provenance claims, not a report of source access undertaken for this draft.
\end{itemize}
\end{document}
